\begin{definitionbox}{Definition (Semigroup)}
\begin{definition}[Semigroup]
\label{def:semigroup}
A \emph{semigroup} is a domain with a closed associative binary operation.
\end{definition}
\end{definitionbox}

\begin{definitionbox}{Definition (Monoid)}
\begin{definition}[Monoid]
\label{def:monoid}
A \emph{monoid} is a semigroup with a two-sided identity element.
\end{definition}
\end{definitionbox}

\begin{definitionbox}{Definition (Commutative Monoid)}
\begin{definition}[Commutative Monoid]
\label{def:commutative-monoid}
A \emph{commutative monoid} is a monoid whose operation is commutative.
\end{definition}
\end{definitionbox}

\begin{definitionbox}{Definition (Group)}
\begin{definition}[Group]
\label{def:group}
A \emph{group} is a monoid in which every element has a two-sided inverse.
\end{definition}
\end{definitionbox}

\begin{definitionbox}{Definition (Abelian Group)}
\begin{definition}[Abelian Group]
\label{def:abelian-group}
An \emph{abelian group} is a group whose operation is commutative.
\end{definition}
\end{definitionbox}

\begin{definitionbox}{Definition (Semiring)}
\begin{definition}[Semiring]
\label{def:semiring}
A \emph{semiring} is a domain with addition and multiplication such that
addition is a commutative monoid, multiplication is a monoid, multiplication
distributes over addition, and the additive identity is absorbing for
multiplication.
\end{definition}
\end{definitionbox}

\begin{definitionbox}{Definition (Ring)}
\begin{definition}[Ring]
\label{def:ring}
A \emph{ring} is a domain whose additive structure is an abelian group, whose
multiplicative structure is a monoid, and whose multiplication distributes
over addition.
\end{definition}
\end{definitionbox}

\begin{definitionbox}{Definition (Commutative Ring)}
\begin{definition}[Commutative Ring]
\label{def:commutative-ring}
A \emph{commutative ring} is a ring whose multiplication is commutative.
\end{definition}
\end{definitionbox}

\begin{definitionbox}{Definition (Integral Domain)}
\begin{definition}[Integral Domain]
\label{def:abstract-integral-domain}
An \emph{integral domain} is a commutative ring with $0\ne1$ and no zero
divisors.
\end{definition}
\end{definitionbox}

\begin{definitionbox}{Definition (Field)}
\begin{definition}[Field]
\label{def:abstract-field}
A \emph{field} is a commutative ring in which every nonzero element has a
multiplicative inverse.
\end{definition}
\end{definitionbox}

\begin{definitionbox}{Definition (Ordered Field)}
\begin{definition}[Ordered Field]
\label{def:abstract-ordered-field}
An \emph{ordered field} is a field equipped with a total order compatible with
addition and multiplication by positive elements.
\end{definition}
\end{definitionbox}

\begin{axiombox}{Axiom (Ordered Field Trichotomy)}
\begin{axiom}[Ordered Field Trichotomy]
\label{ax:ordered-field-trichotomy}
In an ordered field, for all elements $a$ and $b$, exactly one of the following
holds:
\[
a=b,\qquad a<b,\qquad b<a.
\]
\end{axiom}
\end{axiombox}

\begin{remark*}[Interpretation]
Any two elements of an ordered field are comparable, and comparison has no
overlap: equality, being below, and being above are mutually exclusive cases.
\end{remark*}

\NoLocalDependencies

\begin{axiombox}{Axiom (Ordered Field Addition Law)}
\begin{axiom}[Ordered Field Addition Law]
\label{ax:ordered-field-addition-law}
In an ordered field, for all elements $a$, $b$, and $c$,
\[
a<b \quad\Longleftrightarrow\quad a+c<b+c.
\]
\end{axiom}
\end{axiombox}

\begin{remark*}[Interpretation]
Adding the same element to both sides of a strict inequality preserves and
reflects the order.
\end{remark*}

\NoLocalDependencies

\begin{axiombox}{Axiom (Ordered Field Multiplication Law)}
\begin{axiom}[Ordered Field Multiplication Law]
\label{ax:ordered-field-multiplication-law}
In an ordered field, for all elements $a$, $b$, and $c$,
\[
c>0 \Longrightarrow
\bigl(ac<bc \Longleftrightarrow a<b\bigr),
\]
and
\[
c<0 \Longrightarrow
\bigl(ac<bc \Longleftrightarrow b<a\bigr).
\]
Moreover, positive elements are closed under multiplication:
\[
0<a \land 0<b \Longrightarrow 0<ab.
\]
\end{axiom}
\end{axiombox}

\begin{remark*}[Interpretation]
Multiplication by a positive element preserves strict order, while
multiplication by a negative element reverses strict order. The positive
elements are also closed under multiplication.
\end{remark*}

\NoLocalDependencies

\begin{theorembox}{Theorem (Ordered Field Transitivity)}
\begin{theorem}[Ordered Field Transitivity]
\label{thm:ordered-field-transitivity}
In an ordered field, for all elements $a$, $b$, and $c$,
\[
a<b \land b<c \Longrightarrow a<c.
\]
\end{theorem}
\end{theorembox}

\begin{remark*}[Interpretation]
Strict order in an ordered field chains: an element below a second element
which is below a third is below the third.
\end{remark*}

\begin{dependencies}
\begin{itemize}
  \item \hyperref[def:abstract-ordered-field]{Ordered Field}
\end{itemize}
\end{dependencies}

\begin{definitionbox}{Definition (Positive Cone)}
\begin{definition}[Positive cone]
\label{def:positive-cone}
Let $F$ be a field. A subset $P\subset F$ is called a \emph{positive
cone} if it satisfies:
\begin{enumerate}
  \item if $a,b\in P$, then $a+b\in P$;
  \item if $a,b\in P$, then $ab\in P$;
  \item for every $a\in F$, exactly one of the following holds:
  \[
  a\in P,
  \qquad
  a=0,
  \qquad
  -a\in P.
  \]
\end{enumerate}
Define the order relation by
\[
a<b \iff b-a\in P.
\]
\end{definition}
\end{definitionbox}
